\setcounter{chapter}{5}
\chapter{Carathéodory流の測度論}

\paragraph{本章の目標}
これまでのLebesgue流
\[
\text{外測度}=\text{内測度}
\]
という定義を，より抽象的で構造的なCarathéodory流の定義へ結びつける．
さらに，区間上の体積から外測度を作り，
そこから測度を得る一般的な構成法を見る．

\section{動機づけ}

第4章で導入した内測度は，
幾何学的には非常に自然であった．
実際，
\[
\lambda_*(A)
=
\sup\{\lambda^*(K)\mid K \subset A,\ K \text{ はコンパクト}\}
\]
という式は，
「$A$ の中にどれだけ大きなコンパクト集合が入るか」
という直観そのものである．

しかし，この定義には限界がある．
\begin{enumerate}
    \item コンパクト集合という位相的概念に依存している
    \item 一般の集合 $X$ 上では同じ定義ができない
    \item 可測集合全体が $\sigma$-加法族になることを直接示すのは難しい
\end{enumerate}

そこでCarathéodoryは，
内測度を使わず，外測度だけから可測集合を取り出す判定条件を与えた．
この方法の利点は，位相を持たない一般の集合 $X$ に対しても
そのまま適用できることにある．

\section{Carathéodory可測性}

\begin{definition}[Carathéodory可測]
集合 $X$ 上の外測度 $\mu^*$ を取る．
集合 $A \subset X$ がCarathéodory可測であるとは，
任意の $E \subset X$ に対して
\[
\mu^*(E)
=
\mu^*(E \cap A)+\mu^*(E \setminus A)
\]
が成り立つことをいう．
\end{definition}

\begin{remark}
外測度の可算劣加法性から
\[
\mu^*(E)
\le
\mu^*(E \cap A)+\mu^*(E \setminus A)
\]
は常に成り立つ．
したがってCarathéodory条件で本質的なのは逆向きの不等式である．
これは，$A$ で切り分けても外測度が余計に増えないことを表している．
\end{remark}

\begin{definition}
外測度 $\mu^*$ に関するCarathéodory可測集合全体を
\[
\mathcal M_{\mu^*}
\]
と書く．
\end{definition}

\begin{lemma}
\[
\emptyset \in \mathcal M_{\mu^*},\qquad X \in \mathcal M_{\mu^*}
\]
が成り立つ．
\end{lemma}
\begin{proof}
任意の $E \subset X$ に対して
\[
E \cap \emptyset = \emptyset,\qquad E \setminus \emptyset = E
\]
だから
\[
\mu^*(E)=0+\mu^*(E)
\]
である．
したがって $\emptyset \in \mathcal M_{\mu^*}$ である．

また
\[
E \cap X = E,\qquad E \setminus X = \emptyset
\]
だから同様に $X \in \mathcal M_{\mu^*}$ である．
\end{proof}

\begin{lemma}[補集合に関する閉性]
$A \in \mathcal M_{\mu^*}$ ならば $X \setminus A \in \mathcal M_{\mu^*}$ である．
\end{lemma}
\begin{proof}
任意の $E \subset X$ に対して
\[
E \cap (X\setminus A)=E\setminus A,\qquad
E\setminus (X\setminus A)=E\cap A
\]
だから，
$A$ の可測性より
\[
\mu^*(E)
=
\mu^*(E\cap A)+\mu^*(E\setminus A)
=
\mu^*(E\setminus (X\setminus A))+\mu^*(E\cap (X\setminus A)).
\]
したがって $X\setminus A \in \mathcal M_{\mu^*}$ である．
\end{proof}

\begin{lemma}[有限和に関する閉性]
$A,B \in \mathcal M_{\mu^*}$ ならば $A \cup B \in \mathcal M_{\mu^*}$ である．
\end{lemma}
\begin{proof}
任意の $E \subset X$ を取る．
まず $A$ の可測性より
\[
\mu^*(E)=\mu^*(E\cap A)+\mu^*(E\setminus A)
\]
である．
ここで $B$ の可測性を $E\cap A$ と $E\setminus A$ に適用すると
\[
\mu^*(E\cap A)
=
\mu^*(E\cap A\cap B)+\mu^*(E\cap A\setminus B)
\]
および
\[
\mu^*(E\setminus A)
=
\mu^*(E\cap B\setminus A)+\mu^*(E\setminus (A\cup B))
\]
を得る．
したがって
\begin{align*}
\mu^*(E)
=\,&
\mu^*(E\cap A\cap B)
+\mu^*(E\cap A\setminus B) \\
&+\mu^*(E\cap B\setminus A)
+\mu^*(E\setminus (A\cup B)).
\end{align*}
一方，
\[
E\cap (A\cup B)
=
(E\cap A\cap B)\sqcup(E\cap A\setminus B)\sqcup(E\cap B\setminus A)
\]
だから，外測度の可算劣加法性より
\[
\mu^*(E\cap (A\cup B))
\le
\mu^*(E\cap A\cap B)
+\mu^*(E\cap A\setminus B)
+\mu^*(E\cap B\setminus A).
\]
ゆえに
\[
\mu^*(E)
\ge
\mu^*(E\cap (A\cup B))+\mu^*(E\setminus (A\cup B)).
\]
逆向きの不等式は外測度の劣加法性から常に成り立つので，
\[
\mu^*(E)
=
\mu^*(E\cap (A\cup B))+\mu^*(E\setminus (A\cup B))
\]
である．
したがって $A\cup B \in \mathcal M_{\mu^*}$ である．
\end{proof}

\begin{lemma}
$\mathcal M_{\mu^*}$ は有限共通部分と差集合に関して閉じている．
\end{lemma}
\begin{proof}
補集合と有限和に関する閉性から
\[
A\cap B = X\setminus\bigl((X\setminus A)\cup(X\setminus B)\bigr)
\]
および
\[
A\setminus B=A\cap(X\setminus B)
\]
により従う．
\end{proof}

\begin{lemma}[可算和に関する閉性]
$\{A_n\}_{n=1}^{\infty}\subset \mathcal M_{\mu^*}$ ならば
\[
\bigcup_{n=1}^{\infty}A_n \in \mathcal M_{\mu^*}
\]
である．
\end{lemma}
\begin{proof}
まず
\[
B_1:=A_1,\qquad
B_n:=A_n\setminus \bigcup_{k=1}^{n-1}A_k\quad(n\ge 2)
\]
とおく．
前の補題より各 $B_n$ は可測であり，
さらに $\{B_n\}$ は互いに素で
\[
\bigcup_{n=1}^{\infty}B_n=\bigcup_{n=1}^{\infty}A_n
\]
である．
したがって，互いに素な可測集合列の場合を示せば十分である．

そこで $\{B_n\}$ を互いに素な可測集合列とし，
\[
B:=\bigcup_{n=1}^{\infty}B_n
\]
とおく．
任意の $E \subset X$ と $N \in \NN$ に対して，有限和に関する閉性を繰り返し用いると
\[
\mu^*(E)
=
\sum_{n=1}^{N}\mu^*(E\cap B_n)
+\mu^*\left(E\setminus \bigcup_{n=1}^{N}B_n\right)
\]
を得る．
ここで
\[
E\setminus B \subset E\setminus \bigcup_{n=1}^{N}B_n
\]
だから単調性より
\[
\mu^*(E\setminus B)
\le
\mu^*\left(E\setminus \bigcup_{n=1}^{N}B_n\right).
\]
したがって
\[
\mu^*(E)
\ge
\sum_{n=1}^{N}\mu^*(E\cap B_n)+\mu^*(E\setminus B)
\]
が成り立つ．
$N$ について上限を取ると
\[
\mu^*(E)
\ge
\sum_{n=1}^{\infty}\mu^*(E\cap B_n)+\mu^*(E\setminus B)
\]
を得る．

一方，
\[
E\cap B=\bigcup_{n=1}^{\infty}(E\cap B_n)
\]
だから可算劣加法性より
\[
\mu^*(E\cap B)
\le
\sum_{n=1}^{\infty}\mu^*(E\cap B_n).
\]
よって
\[
\mu^*(E)
\ge
\mu^*(E\cap B)+\mu^*(E\setminus B).
\]
逆向きの不等式は外測度の劣加法性から常に成り立つので，
\[
\mu^*(E)
=
\mu^*(E\cap B)+\mu^*(E\setminus B)
\]
である．
したがって $B \in \mathcal M_{\mu^*}$ である．
\end{proof}

\begin{theorem}[Carathéodory]
$\mathcal M_{\mu^*}$ は $X$ 上の $\sigma$-加法族である．
さらに，制限
\[
\mu:=\mu^*|_{\mathcal M_{\mu^*}}
\]
は $(X,\mathcal M_{\mu^*})$ 上の測度であり，
零集合の部分集合も再び可測である．
\end{theorem}
\begin{proof}
前の補題たちにより，$\mathcal M_{\mu^*}$ は $X$ を含み，
補集合と可算和に関して閉じている．
したがって $\sigma$-加法族である．

次に $\mu$ が測度であることを示す．
$\{A_n\}_{n=1}^{\infty}\subset \mathcal M_{\mu^*}$ を互いに素な列とし，
\[
A:=\bigcup_{n=1}^{\infty}A_n
\]
とおく．
上で示した可算和の証明の途中式を $E=A$ に適用すると，
任意の $N$ に対して
\[
\mu(A)
=
\sum_{n=1}^{N}\mu(A_n)
+\mu\left(\bigcup_{n=N+1}^{\infty}A_n\right)
\]
である．
特に
\[
\mu(A)\ge \sum_{n=1}^{N}\mu(A_n)
\]
がすべての $N$ に対して成り立つから
\[
\mu(A)\ge \sum_{n=1}^{\infty}\mu(A_n)
\]
である．
逆向きの不等式は外測度の可算劣加法性から
\[
\mu(A)=\mu^*(A)\le \sum_{n=1}^{\infty}\mu^*(A_n)=\sum_{n=1}^{\infty}\mu(A_n)
\]
である．
したがって
\[
\mu(A)=\sum_{n=1}^{\infty}\mu(A_n)
\]
を得る．
また $\mu(\emptyset)=\mu^*(\emptyset)=0$ だから，$\mu$ は測度である．

最後に，$N \in \mathcal M_{\mu^*}$ で $\mu(N)=0$ とし，
$A \subset N$ を任意に取る．
任意の $E \subset X$ に対して，
外測度の劣加法性より
\[
\mu^*(E)\le \mu^*(E\cap A)+\mu^*(E\setminus A)
\]
である．
一方，
\[
E\cap A \subset N
\]
だから単調性より
\[
\mu^*(E\cap A)\le \mu^*(N)=\mu(N)=0
\]
であり，
\[
\mu^*(E\cap A)=0
\]
を得る．
さらに
\[
E\setminus A \subset E
\]
より
\[
\mu^*(E\setminus A)\le \mu^*(E)
\]
である．
したがって
\[
\mu^*(E\cap A)+\mu^*(E\setminus A)\le \mu^*(E)
\]
となる．
両側の不等式を合わせると
\[
\mu^*(E)=\mu^*(E\cap A)+\mu^*(E\setminus A)
\]
である．
よって $A \in \mathcal M_{\mu^*}$ である．
\end{proof}

\section{Lebesgue流との関係}

\begin{proposition}[閉区間はCarathéodory可測]
有界閉区間
\[
I=[a_1,b_1]\times\cdots\times[a_d,b_d]
\]
は $\lambda^*$ に関してCarathéodory可測である．
\end{proposition}
\begin{proof}
任意の $E \subset \RR^d$ と $\eps>0$ を取る．
第3章の開区間被覆の定理より，
$E$ を覆う有界開区間列 $\{U_n\}_{n=1}^{\infty}$ で
\[
\sum_{n=1}^{\infty}|U_n|<\lambda^*(E)+\eps
\]
を満たすものが取れる．

各 $n$ について
\[
U_n=(\alpha_1^{(n)},\beta_1^{(n)})\times\cdots\times(\alpha_d^{(n)},\beta_d^{(n)})
\]
と書く．
各座標方向で区間
\[
(\alpha_i^{(n)},a_i),\qquad
(a_i,b_i),\qquad
(b_i,\beta_i^{(n)})
\]
のうち空でないものだけを取り出すと，
$U_n$ は有限個の互いに素な有界開区間に分割される．
そのうち $I$ と交わる部分の合併は $U_n\cap I$ であり，
残りの合併は $U_n\setminus I$ である．
したがって $U_n\cap I$ と $U_n\setminus I$ はそれぞれ有限個の有界開区間の合併であり，
その体積の総和は
\[
|U_n|
\]
に等しい．

ゆえに $E\cap I$ と $E\setminus I$ は，それぞれこれらの開区間族で覆われるから
\[
\lambda^*(E\cap I)+\lambda^*(E\setminus I)
\le
\sum_{n=1}^{\infty}|U_n|
<
\lambda^*(E)+\eps
\]
である．
$\eps>0$ は任意だから
\[
\lambda^*(E\cap I)+\lambda^*(E\setminus I)\le \lambda^*(E)
\]
を得る．
逆向きの不等式は外測度の劣加法性から常に成り立つので
\[
\lambda^*(E)
=
\lambda^*(E\cap I)+\lambda^*(E\setminus I)
\]
である．
したがって $I$ はCarathéodory可測である．
\end{proof}

\begin{proposition}[開集合の可測性]
任意の開集合 $G \subset \RR^d$ は $\lambda^*$ に関してCarathéodory可測である．
\end{proposition}
\begin{proof}
まず第5章で示したように，座標超平面片は外測度 $0$ をもつ．
したがってCarathéodoryの定理より，
それらの部分集合はすべてCarathéodory可測である．

次に任意の有界開区間
\[
J=(a_1,b_1)\times\cdots\times(a_d,b_d)
\]
を取る．
同じ端点をもつ閉区間
\[
\overline J=[a_1,b_1]\times\cdots\times[a_d,b_d]
\]
との差
\[
\overline J\setminus J
\]
は有限個の座標超平面片の合併に含まれるからCarathéodory可測である．
閉区間 $\overline J$ も前命題よりCarathéodory可測だから，
\[
J=\overline J\setminus(\overline J\setminus J)
\]
もCarathéodory可測である．

次に，端点がすべて有理数である有界開区間全体の族を $\mathcal J$ とする．
$\mathcal J$ は $\QQ^{2d}$ と1対1に対応するので可算集合である．
開集合 $G$ を任意に取る．
各 $x \in G$ に対して，$G$ の開性より $r>0$ で
\[
[x_1-r,x_1+r]\times\cdots\times[x_d-r,x_d+r]\subset G
\]
を満たすものが取れる．
有理数の稠密性により，各 $i$ について
\[
x_i \in (p_i,q_i)\subset [x_i-r,x_i+r]
\]
を満たす有理数 $p_i,q_i$ が取れる．
すると
\[
J_x:=(p_1,q_1)\times\cdots\times(p_d,q_d)
\]
は $\mathcal J$ の元で
\[
x\in J_x\subset G
\]
を満たす．
よって
\[
G=\bigcup\{J \in \mathcal J\mid J\subset G\}
\]
である．
右辺は可算個のCarathéodory可測集合の合併だから，Carathéodory可測である．
\end{proof}

\begin{corollary}[Borel集合]
すべてのBorel集合は $\lambda^*$ に関してCarathéodory可測である．
\end{corollary}
\begin{proof}
Carathéodory可測集合全体はCarathéodoryの定理により $\sigma$-加法族であり，
前命題よりすべての開集合を含む．
したがって，開集合が生成するBorel集合族全体を含む．
\end{proof}

\begin{proposition}[開集合の内正則性]
任意の開集合 $G \subset \RR^d$ に対して
\[
\lambda^*(G)
=
\sup\{\lambda^*(K)\mid K \subset G,\ K \text{ はコンパクト}\}
\]
が成り立つ．
\end{proposition}
\begin{proof}
前命題の証明で用いた可算集合
\[
\mathcal J_G:=\{J \in \mathcal J\mid \overline J\subset G\}
\]
を列
\[
\mathcal J_G=\{J_1,J_2,\dots\}
\]
と書く．
すると
\[
G=\bigcup_{n=1}^{\infty}J_n
\]
である．
各 $N$ に対して
\[
U_N:=\bigcup_{n=1}^{N}J_n,\qquad
K_N:=\bigcup_{n=1}^{N}\overline{J_n}
\]
とおく．
$K_N$ は有限個のコンパクト集合の合併だからコンパクトであり，
\[
K_N\subset G
\]
である．

また各 $\overline{J_n}\setminus J_n$ は有限個の座標超平面片の合併に含まれるから外測度 $0$ をもつ．
したがって
\[
\lambda^*(\overline{J_n})=\lambda^*(J_n)
\]
である．
さらに有限個の互いに素な有界開区間の合併に対しては，
Carathéodoryの定理から得られた測度の有限加法性により
\[
\lambda^*(K_N)=\lambda^*(U_N)
\]
が成り立つ．

いま
\[
U_1\subset U_2\subset \cdots,\qquad
G=\bigcup_{N=1}^{\infty}U_N
\]
だから，Carathéodoryの定理から得られた測度の下からの連続性より
\[
\lambda^*(G)=\lim_{N\to\infty}\lambda^*(U_N)=\lim_{N\to\infty}\lambda^*(K_N).
\]
したがって
\[
\lambda^*(G)\le \sup\{\lambda^*(K)\mid K\subset G,\ K\text{ はコンパクト}\}
\]
である．
逆向きの不等式は外測度の単調性から明らかである．
\end{proof}

\begin{theorem}[Carathéodory可測ならLebesgue可測]
$A \subset \RR^d$ が $\lambda^*$ に関してCarathéodory可測ならば
\[
\lambda^*(A)=\lambda_*(A)
\]
が成り立つ．
\end{theorem}
\begin{proof}
$A$ がCarathéodory可測であるとする．
第3章の外正則性より，任意の $\eps>0$ に対して
$A \subset G$ を満たす開集合 $G$ で
\[
\lambda^*(G)<\lambda^*(A)+\eps
\]
を満たすものが取れる．
開集合 $G$ はCarathéodory可測だから，
Carathéodoryの定理から得られる測度を再び $\lambda^*$ と書くことにすれば
\[
\lambda^*(G\setminus A)=\lambda^*(G)-\lambda^*(A)<\eps
\]
である．

さらに前命題より，$G$ の中にコンパクト集合 $K$ で
\[
\lambda^*(G)-\lambda^*(K)<\eps
\]
を満たすものが取れる．
また第3章の外正則性を $G\setminus A$ に適用すると，
$G\setminus A \subset O$ を満たす開集合 $O$ で
\[
\lambda^*(O)<\lambda^*(G\setminus A)+\eps<2\eps
\]
となるものが取れる．

\[
F:=K\setminus O
\]
とおくと，$F$ はコンパクトであり，
さらに
\[
F\subset G\setminus O\subset A
\]
である．
したがって
\[
\lambda_*(A)\ge \lambda^*(F)
\]
である．
一方，$F\subset K$ だから
\[
K=(K\cap O)\sqcup F
\]
であり，加法性より
\[
\lambda^*(K)=\lambda^*(K\cap O)+\lambda^*(F)\le \lambda^*(O)+\lambda^*(F)
\]
となる．
したがって
\[
\lambda^*(F)\ge \lambda^*(K)-\lambda^*(O).
\]
以上をまとめると
\[
\lambda_*(A)\ge \lambda^*(K)-\lambda^*(O)>\lambda^*(G)-3\eps>\lambda^*(A)-3\eps.
\]
もともと常に
\[
\lambda_*(A)\le \lambda^*(A)
\]
だから，$\eps>0$ は任意であることより
\[
\lambda_*(A)=\lambda^*(A)
\]
を得る．
\end{proof}

\begin{theorem}[Lebesgue可測ならCarathéodory可測]
$A \subset \RR^d$ が
\[
\lambda^*(A)=\lambda_*(A)
\]
を満たすならば，$A$ は $\lambda^*$ に関してCarathéodory可測である．
\end{theorem}
\begin{proof}
$A$ がLebesgue可測であるとする．
任意の $\eps>0$ に対して，
第3章の外正則性より $A \subset G$ を満たす開集合 $G$ で
\[
\lambda^*(G)<\lambda^*(A)+\eps
\]
となるものが取れる．
また内測度の定義より，コンパクト集合 $K \subset A$ で
\[
\lambda^*(K)>\lambda_*(A)-\eps=\lambda^*(A)-\eps
\]
を満たすものが取れる．

開集合 $G$ とコンパクト集合 $K$ はいずれもCarathéodory可測だから，
\[
G\setminus K
\]
もCarathéodory可測であり，
\[
\lambda^*(G\setminus K)=\lambda^*(G)-\lambda^*(K)<2\eps
\]
である．

任意の $E \subset \RR^d$ を取る．
$K$ のCarathéodory可測性より
\[
\lambda^*(E)=\lambda^*(E\cap K)+\lambda^*(E\setminus K)
\]
である．
ここで
\[
E\cap A \subset (E\cap K)\cup(A\setminus K)
\]
だから劣加法性より
\[
\lambda^*(E\cap A)\le \lambda^*(E\cap K)+\lambda^*(A\setminus K).
\]
また
\[
A\setminus K \subset G\setminus K
\]
だから
\[
\lambda^*(A\setminus K)\le \lambda^*(G\setminus K)<2\eps.
\]
したがって
\[
\lambda^*(E\cap K)>\lambda^*(E\cap A)-2\eps.
\]
さらに
\[
E\setminus A \subset E\setminus K
\]
だから単調性より
\[
\lambda^*(E\setminus K)\ge \lambda^*(E\setminus A).
\]
以上を合わせると
\[
\lambda^*(E)
>
\lambda^*(E\cap A)+\lambda^*(E\setminus A)-2\eps
\]
を得る．
$\eps>0$ は任意だから
\[
\lambda^*(E)\ge \lambda^*(E\cap A)+\lambda^*(E\setminus A)
\]
である．
逆向きの不等式は外測度の劣加法性から常に成り立つので，
\[
\lambda^*(E)
=
\lambda^*(E\cap A)+\lambda^*(E\setminus A)
\]
となる．
したがって $A$ はCarathéodory可測である．
\end{proof}

\begin{corollary}[2つの定義の一致]
$\RR^d$ 上では，Lebesgue流の可測性
\[
\lambda^*(A)=\lambda_*(A)
\]
と，Carathéodory流の可測性とは同値である．
\end{corollary}
\begin{proof}
直前の2つの定理を合わせればよい．
\end{proof}

\begin{corollary}[第5章の基本定理]
\[
\mathcal L^d
=
\mathcal M_{\lambda^*}
\]
である．
したがって $\mathcal L^d$ は $\RR^d$ 上の $\sigma$-加法族であり，
\[
\lambda:\mathcal L^d\to[0,\infty]
\]
は測度である．
\end{corollary}
\begin{proof}
前系とCarathéodoryの定理より直ちに従う．
\end{proof}

\begin{remark}
ここで初めて，
第4章の内測度が「補助的」であった理由がはっきりする．
Lebesgue測度のような具体的な幾何学の場面では
内測度は非常に有用であるが，
測度論そのものを組み立てるときの主役は
外測度とCarathéodory条件である．
\end{remark}

\begin{theorem}[有界集合に対する内測度の別定義]
$A \subset \RR^d$ を有界集合とし，
$A \subset I$ を満たす有界閉区間 $I$ を1つ取る．
このとき
\[
\lambda_*(A)=|I|-\lambda^*(I\setminus A)
\]
が成り立つ．
\end{theorem}
\begin{proof}
まず
\[
\alpha(A):=|I|-\lambda^*(I\setminus A)
\]
とおく．

$K \subset A$ を任意のコンパクト集合とする．
第5章より $K$ はLebesgue可測であり，$I$ もコンパクトだからLebesgue可測である．
また
\[
I\setminus K
\]
は可測集合の差だからLebesgue可測であり，
\[
I=K\sqcup(I\setminus K)
\]
である．
したがって測度の加法性より
\[
|I|=\lambda(I)=\lambda(K)+\lambda(I\setminus K)=\lambda^*(K)+\lambda(I\setminus K)
\]
である．
さらに
\[
I\setminus A \subset I\setminus K
\]
だから単調性より
\[
\lambda^*(I\setminus A)\le \lambda^*(I\setminus K)=\lambda(I\setminus K)
\]
である．
よって
\[
\lambda^*(K)\le |I|-\lambda^*(I\setminus A)=\alpha(A)
\]
を得る．
$K$ の取り方について上限を取れば
\[
\lambda_*(A)\le \alpha(A)
\]
である．

逆向きの不等式を示す．
$\eps>0$ を任意に取る．
第3章の外正則性より，
$I\setminus A \subset G$ を満たす開集合 $G$ で
\[
\lambda^*(G)<\lambda^*(I\setminus A)+\eps
\]
となるものが取れる．
ここで
\[
K:=I\setminus G
\]
とおくと，$K$ はコンパクトであり，しかも
\[
K\subset A
\]
である．
また $G$ は開集合だから第6章前半で示したようにLebesgue可測であり，
$I\cap G$ も可測である．
さらに
\[
I=K\sqcup(I\cap G)
\]
だから
\[
\lambda(K)=\lambda(I)-\lambda(I\cap G)=|I|-\lambda(I\cap G).
\]
いま
\[
I\cap G\subset G
\]
だから単調性より
\[
\lambda(I\cap G)\le \lambda(G)=\lambda^*(G)
\]
である．
したがって
\[
\lambda^*(K)=\lambda(K)\ge |I|-\lambda^*(G)> |I|-\lambda^*(I\setminus A)-\eps=\alpha(A)-\eps.
\]
ゆえに
\[
\lambda_*(A)\ge \lambda^*(K)>\alpha(A)-\eps.
\]
$\eps>0$ は任意だから
\[
\lambda_*(A)\ge \alpha(A)
\]
である．

以上より
\[
\lambda_*(A)=\alpha(A)=|I|-\lambda^*(I\setminus A)
\]
を得る．
\end{proof}

\begin{corollary}[区間の取り方によらない]
$A$ を有界集合とし，
\[
A\subset I,\qquad A\subset J
\]
を満たす有界閉区間 $I,J$ を取る．
このとき
\[
|I|-\lambda^*(I\setminus A)=|J|-\lambda^*(J\setminus A)
\]
が成り立つ．
\end{corollary}
\begin{proof}
どちらも直前の定理により $\lambda_*(A)$ に等しい．
\end{proof}

\section{Borel集合と非可測集合}

\begin{theorem}[Borel集合]
すべてのBorel集合はLebesgue可測である．
\end{theorem}
\begin{proof}
前節で，すべてのBorel集合は $\lambda^*$ に関してCarathéodory可測であることを示した．
またLebesgue可測性とCarathéodory可測性は一致するから，主張が従う．
\end{proof}

\begin{theorem}[Borel測度の完備化]
任意のLebesgue可測集合 $A$ は，
あるBorel集合 $B$ とあるLebesgue零集合 $N$ を用いて
\[
A=B\triangle N
\]
と書ける．
したがって，Lebesgue測度はBorel測度の完備化として得られる．
\end{theorem}
\begin{proof}
$A$ をLebesgue可測集合とする．
各 $n\in\NN$ に対して，第3章の外正則性より
$A \subset G_n$ を満たす開集合 $G_n$ で
\[
\lambda(G_n\setminus A)<\frac1n
\]
を満たすものが取れる．
ここで
\[
B:=\bigcap_{n=1}^{\infty}G_n
\]
とおく．
$B$ は開集合の可算共通部分だからBorel集合であり，
しかも $A \subset B$ である．

いま
\[
B\setminus A \subset G_n\setminus A
\]
がすべての $n$ について成り立つから，単調性より
\[
0\le \lambda(B\setminus A)\le \lambda(G_n\setminus A)<\frac1n
\]
である．
したがって
\[
\lambda(B\setminus A)=0.
\]
$N:=B\setminus A$ とおけば
\[
A=B\setminus N=B\triangle N
\]
であり，しかも $N$ は零集合である．
\end{proof}

\begin{theorem}[Vitali]
$[0,1]$ にはLebesgue非可測集合が存在する．
\end{theorem}
\begin{proof}
$x-y \in \QQ$ によって $[0,1]$ 上に同値関係を定める．
選択公理を用いて，各同値類からちょうど1点ずつ選んだ集合を $V \subset [0,1]$ とする．

有理数
\[
q\in [-1,1]\cap\QQ
\]
に対して
\[
V+q:=\{v+q\mid v\in V\}
\]
を考える．
まず，$q_1\ne q_2$ なら $(V+q_1)\cap(V+q_2)=\emptyset$ である．
実際，もし
\[
x=v_1+q_1=v_2+q_2
\]
となる $v_1,v_2\in V$ があれば
\[
v_1-v_2=q_2-q_1\in\QQ
\]
だから $v_1,v_2$ は同じ同値類に属する．
しかし $V$ は各同値類から1点しか選んでいないので $v_1=v_2$ であり，
従って $q_1=q_2$ となる．
これは仮定に反する．

次に
\[
[0,1]\subset \bigcup_{q\in[-1,1]\cap\QQ}(V+q)\subset [-1,2]
\]
が成り立つ．
左側の包含は，任意の $x\in[0,1]$ に対して，$x$ と同じ同値類に属する代表元 $v\in V$ が存在し，
\[
x-v\in\QQ\cap[-1,1]
\]
となることから従う．
右側の包含は $V\subset [0,1]$ だから明らかである．

もし $V$ がLebesgue可測なら，各 $V+q$ も平行移動不変性によりLebesgue可測で，
\[
\lambda_1(V+q)=\lambda_1(V)
\qquad
(q\in[-1,1]\cap\QQ)
\]
である．
すると可算加法性より
\[
\lambda_1\left(\bigcup_{q\in[-1,1]\cap\QQ}(V+q)\right)
=
\sum_{q\in[-1,1]\cap\QQ}\lambda_1(V).
\]

\begin{enumerate}
    \item $\lambda_1(V)>0$ なら右辺は $\infty$ になるが，
    左辺は $[-1,2]$ に含まれるので高々 $3$ である．
    これは矛盾である．
    \item $\lambda_1(V)=0$ なら右辺は $0$ になるが，
    左辺は $[0,1]$ を含むので少なくとも $1$ である．
    これも矛盾である．
\end{enumerate}

したがって $V$ はLebesgue可測ではない．
\end{proof}

\begin{remark}
この定理は，「すべての部分集合に面積を与える」ことはできないことを示している．
Lebesgue測度は非常に広いクラスの集合を扱えるが，
それでもなお全能ではない．
\end{remark}

\begin{remark}[任意加法性は強すぎる]
Lebesgue測度は可算加法的であるが，
任意個の互いに素な集合族に対してまで加法性を要求することはできない．
実際，すべての1点集合は測度 $0$ であるから，
もし任意加法性
\[
\lambda\left(\bigsqcup_{\alpha \in \Gamma} A_\alpha\right)
=
\sum_{\alpha \in \Gamma} \lambda(A_\alpha)
\]
が常に成り立つとする
（右辺は有限部分和の上限として解釈する）．
すると
\[
[0,1] = \bigsqcup_{x \in [0,1]} \{x\}
\]
より
\[
\lambda_1([0,1]) = \sum_{x \in [0,1]} 0 = 0
\]
となってしまう．
これは明らかに矛盾である．
したがって，面積を正しくモデル化するには
「可算加法性」がちょうどよい強さである．
\end{remark}

\section{一般の測度の構成}

\begin{definition}[有限加法族]
集合 $X$ の部分集合族 $\mathcal A \subset \mathcal P(X)$ が
$X$ 上の有限加法族であるとは，次の3条件を満たすことをいう．
\begin{enumerate}
    \item $\emptyset \in \mathcal A$
    \item $A,B \in \mathcal A$ ならば $A \cup B \in \mathcal A$
    \item $A,B \in \mathcal A$ ならば $A \setminus B \in \mathcal A$
\end{enumerate}
\end{definition}

\begin{definition}[半加法族]
集合 $X$ の部分集合族 $\mathcal S \subset \mathcal P(X)$ が
$X$ 上の半加法族であるとは，次の3条件を満たすことをいう．
\begin{enumerate}
    \item $\emptyset \in \mathcal S$
    \item $A,B \in \mathcal S$ ならば $A \cap B \in \mathcal S$
    \item $A,B \in \mathcal S$ ならば，ある有限個の互いに素な
    $C_1,\dots,C_N \in \mathcal S$ が存在して
    \[
    A \setminus B = \bigsqcup_{j=1}^N C_j
    \]
    と書ける
\end{enumerate}
\end{definition}

\begin{definition}[半加法族が生成する有限加法族]
半加法族 $\mathcal S$ に対し，
\[
\mathcal A(\mathcal S)
:=
\left\{
\bigsqcup_{j=1}^N A_j
\;\middle|\;
N \in \NN,\ A_j \in \mathcal S,\ A_i \cap A_j = \emptyset\ (i \ne j)
\right\}
\cup \{\emptyset\}
\]
と定める．
\end{definition}

\begin{example}[半開区間の族]
\[
\mathcal S
:=
\left\{
\prod_{k=1}^d [a_k,b_k)
\;\middle|\;
a_k \le b_k
\right\}
\]
は $\RR^d$ 上の半加法族である．
実際，共通部分は再び半開区間であり，
差集合は有限個の半開区間に分割できる．
第2章の基本集合は，ちょうどこの $\mathcal S$ から作った
$\mathcal A(\mathcal S)$ にほかならない．
\end{example}

\begin{lemma}
$A \in \mathcal S$，$E \in \mathcal A(\mathcal S)$ とする．
このとき $A \setminus E$ は有限個の互いに素な
$\mathcal S$ の元の合併として書ける．
\end{lemma}
\begin{proof}
$E$ の表示に現れる部分集合の個数に関する帰納法で示す．

$E=\emptyset$ のときは
\[
A \setminus E = A \in \mathcal S
\]
であり明らかである．

次に
\[
E=\bigsqcup_{j=1}^N E_j
\qquad
(E_j \in \mathcal S)
\]
で $N \ge 1$ とする．
\[
F:=\bigsqcup_{j=1}^{N-1}E_j
\]
とおくと，帰納法の仮定より
\[
A \setminus F=\bigsqcup_{i=1}^M B_i
\qquad
(B_i \in \mathcal S)
\]
と書ける．
ここで
\[
A \setminus E=(A \setminus F)\setminus E_N
=
\bigsqcup_{i=1}^M (B_i \setminus E_N)
\]
である．
各 $B_i \setminus E_N$ は半加法族の定義より
有限個の互いに素な $\mathcal S$ の元の合併として書ける．
しかも異なる $i$ に属する部分集合は，
それぞれ互いに素な $B_i$ の中に入っているから互いに素である．
したがって $A \setminus E$ も有限個の互いに素な
$\mathcal S$ の元の合併として書ける．
\end{proof}

\begin{proposition}
$\mathcal A(\mathcal S)$ は $X$ 上の有限加法族である．
\end{proposition}
\begin{proof}
\begin{enumerate}
    \item 定義より $\emptyset \in \mathcal A(\mathcal S)$ である．

    \item 差集合に関する閉性を示す．
    \[
    E=\bigsqcup_{i=1}^m A_i,
    \qquad
    F \in \mathcal A(\mathcal S)
    \]
    とする．
    先の補題より，各 $i$ について
    \[
    A_i \setminus F
    \]
    は有限個の互いに素な $\mathcal S$ の元の合併として書ける．
    しかも $A_i$ たちは互いに素だから，
    それらの差集合
    \[
    A_i \setminus F
    \]
    も互いに素である．
    よって
    \[
    E \setminus F
    =
    \bigsqcup_{i=1}^m (A_i \setminus F)
    \]
    は有限個の互いに素な $\mathcal S$ の元の合併であり，
    $E \setminus F \in \mathcal A(\mathcal S)$ である．

    \item 和集合に関する閉性を示す．
    $E,F \in \mathcal A(\mathcal S)$ とすると
    \[
    E \cup F = E \sqcup (F \setminus E)
    \]
    である．
    右辺の $F \setminus E$ はすでに示した差集合に関する閉性より
    $\mathcal A(\mathcal S)$ に属する．
    したがって $E \cup F$ も
    $\mathcal A(\mathcal S)$ に属する．
\end{enumerate}
\end{proof}

\begin{definition}[半加法族上の前測度]
半加法族 $\mathcal S$ 上の写像
\[
\mu_0:\mathcal S \to [0,\infty]
\]
が前測度であるとは，次を満たすことをいう．
\begin{enumerate}
    \item $\mu_0(\emptyset)=0$
    \item 互いに素な列 $\{A_n\}_{n=1}^{\infty}\subset \mathcal S$ で
    \[
    A=\bigsqcup_{n=1}^{\infty}A_n \in \mathcal S
    \]
    となるものに対して
    \[
    \mu_0(A)=\sum_{n=1}^{\infty}\mu_0(A_n)
    \]
    が成り立つ
\end{enumerate}
\end{definition}

\begin{remark}
この定義から有限加法性も従う．
実際，
\[
A=\bigsqcup_{n=1}^N A_n
\]
ならば $A_{N+1}=A_{N+2}=\cdots=\emptyset$ とおいて
可算加法性を適用すればよい．
\end{remark}

\begin{definition}[前測度]
有限加法族 $\mathcal A$ 上の写像
\[
\mu_0:\mathcal A\to[0,\infty]
\]
が前測度であるとは，次を満たすことをいう．
\begin{enumerate}
    \item $\mu_0(\emptyset)=0$
    \item $\mu_0$ は有限加法的である
    \item 互いに素な列 $\{A_n\}_{n=1}^{\infty}\subset \mathcal A$ で
    \[
    \bigcup_{n=1}^{\infty}A_n\in\mathcal A
    \]
    となるものに対して
    \[
    \mu_0\left(\bigcup_{n=1}^{\infty}A_n\right)
    =
    \sum_{n=1}^{\infty}\mu_0(A_n)
    \]
    が成り立つ
\end{enumerate}
\end{definition}

\begin{remark}
単なる有限加法性だけでは一般に測度は構成できない．
可算個の分割と両立するための整合性が必要であり，
それを形式化したものが前測度である．
\end{remark}

\begin{proposition}[半加法族上の前測度の延長]
$\mathcal S$ を半加法族，$\mu_0$ をその上の前測度とする．
各
\[
E=\bigsqcup_{j=1}^N A_j \in \mathcal A(\mathcal S)
\qquad
(A_j \in \mathcal S)
\]
に対して
\[
\overline{\mu}_0(E):=\sum_{j=1}^N \mu_0(A_j)
\]
と定めると，これは well-defined であり，
$\mathcal A(\mathcal S)$ 上の前測度になる．
\end{proposition}
\begin{proof}
まず well-defined であることを示す．
\[
E=\bigsqcup_{i=1}^m A_i=\bigsqcup_{j=1}^n B_j
\qquad
(A_i,B_j \in \mathcal S)
\]
とする．
各 $i$ について
\[
A_i=\bigsqcup_{j=1}^n (A_i \cap B_j)
\]
であり，$A_i \cap B_j \in \mathcal S$ で互いに素である．
したがって半加法族上の前測度の有限加法性より
\[
\mu_0(A_i)=\sum_{j=1}^n \mu_0(A_i \cap B_j)
\]
である．
よって
\[
\sum_{i=1}^m \mu_0(A_i)
=
\sum_{i=1}^m \sum_{j=1}^n \mu_0(A_i \cap B_j).
\]
同様に，各 $j$ について
\[
B_j=\bigsqcup_{i=1}^m (A_i \cap B_j)
\]
だから
\[
\mu_0(B_j)=\sum_{i=1}^m \mu_0(A_i \cap B_j)
\]
であり，
\[
\sum_{j=1}^n \mu_0(B_j)
=
\sum_{j=1}^n \sum_{i=1}^m \mu_0(A_i \cap B_j).
\]
右辺はいずれも同じ有限和だから
\[
\sum_{i=1}^m \mu_0(A_i)=\sum_{j=1}^n \mu_0(B_j)
\]
となり，$\overline{\mu}_0(E)$ は表示に依らない．

次に，$\overline{\mu}_0$ が $\mathcal A(\mathcal S)$ 上の前測度であることを示す．

\begin{enumerate}
    \item $\overline{\mu}_0(\emptyset)=0$ は定義から明らかである．

    \item 有限加法性を示す．
    $E,F \in \mathcal A(\mathcal S)$ が互いに素とし，
    \[
    E=\bigsqcup_{i=1}^m A_i,
    \qquad
    F=\bigsqcup_{j=1}^n B_j
    \]
    と書く．
    $E$ と $F$ は互いに素であるから，
    すべての $A_i$ と $B_j$ も互いに素である．
    よって
    \[
    E \sqcup F
    =
    \left(\bigsqcup_{i=1}^m A_i\right)
    \sqcup
    \left(\bigsqcup_{j=1}^n B_j\right)
    \]
    は $\mathcal A(\mathcal S)$ の表示であり，
    \[
    \overline{\mu}_0(E \sqcup F)
    =
    \sum_{i=1}^m \mu_0(A_i)+\sum_{j=1}^n \mu_0(B_j)
    =
    \overline{\mu}_0(E)+\overline{\mu}_0(F)
    \]
    である．

    \item 可算加法性を示す．
    互いに素な列 $\{E_n\}_{n=1}^{\infty}\subset \mathcal A(\mathcal S)$ で
    \[
    E=\bigsqcup_{n=1}^{\infty}E_n \in \mathcal A(\mathcal S)
    \]
    となるものを取る．
    \[
    E=\bigsqcup_{i=1}^m A_i
    \qquad
    (A_i \in \mathcal S)
    \]
    と書く．
    各 $n,i$ に対して
    \[
    F_{n,i}:=E_n \cap A_i
    \]
    とおくと，$F_{n,i}\in\mathcal A(\mathcal S)$ であり，
    $n$ を動かすと互いに素で
    \[
    A_i=\bigsqcup_{n=1}^{\infty}F_{n,i}
    \]
    が成り立つ．

    ここで各 $F_{n,i}$ を
    \[
    F_{n,i}=\bigsqcup_{k=1}^{r_{n,i}} C_{n,i,k}
    \qquad
    (C_{n,i,k}\in\mathcal S)
    \]
    と書けば，
    \[
    A_i
    =
    \bigsqcup_{n=1}^{\infty}\bigsqcup_{k=1}^{r_{n,i}} C_{n,i,k}
    \]
    である．
    したがって半加法族上の前測度の可算加法性より
    \[
    \mu_0(A_i)
    =
    \sum_{n=1}^{\infty}\sum_{k=1}^{r_{n,i}}\mu_0(C_{n,i,k})
    =
    \sum_{n=1}^{\infty}\overline{\mu}_0(F_{n,i})
    \]
    を得る．
    ゆえに
    \begin{align*}
    \overline{\mu}_0(E)
    &=
    \sum_{i=1}^m \mu_0(A_i) \\
    &=
    \sum_{i=1}^m \sum_{n=1}^{\infty}\overline{\mu}_0(F_{n,i}) \\
    &=
    \sum_{n=1}^{\infty}\sum_{i=1}^m \overline{\mu}_0(F_{n,i}) \\
    &=
    \sum_{n=1}^{\infty}\overline{\mu}_0(E_n).
    \end{align*}
    最後の等号は
    \[
    E_n=\bigsqcup_{i=1}^m F_{n,i}
    \]
    に $\overline{\mu}_0$ の有限加法性を適用したものである．
\end{enumerate}
\end{proof}

\begin{definition}[外測度の生成]
有限加法族 $\mathcal A$ 上の前測度 $\mu_0$ から，
$X$ の任意の部分集合 $E \subset X$ に対して
\[
\mu^*(E)
:=
\inf\left\{
\sum_{n=1}^{\infty}\mu_0(A_n)
\;\middle|\;
E \subset \bigcup_{n=1}^{\infty}A_n,\ 
A_n \in \mathcal A
\right\}
\]
と定める．
\end{definition}

\begin{proposition}
上で定めた $\mu^*$ は $X$ 上の外測度である．
\end{proposition}
\begin{proof}
\begin{enumerate}
    \item $\mu^*(\emptyset)=0$：
    まず $\mu^*(\emptyset)\ge 0$ である．
    一方，
    \[
    \emptyset \subset \emptyset
    \]
    かつ $\emptyset \in \mathcal A$，$\mu_0(\emptyset)=0$ だから
    \[
    \mu^*(\emptyset)\le 0.
    \]
    よって $\mu^*(\emptyset)=0$ である．

    \item 単調性：
    $E \subset F$ とする．
    $F$ の任意の可算被覆はそのまま $E$ の可算被覆でもあるから，
    \[
    \mu^*(E)\le \mu^*(F)
    \]
    である．

    \item 可算劣加法性：
    集合列 $\{E_n\}_{n=1}^{\infty}$ を取る．
    右辺が $\infty$ なら不等式は自明であるから有限とし，
    $\eps>0$ を任意に取る．
    各 $n$ に対して，$\mu^*(E_n)$ の定義より
    \[
    E_n\subset \bigcup_{k=1}^{\infty}A_{n,k},
    \qquad
    \sum_{k=1}^{\infty}\mu_0(A_{n,k})
    <
    \mu^*(E_n)+\frac{\eps}{2^n}
    \]
    となる $A_{n,k}\in\mathcal A$ が取れる．
    すると
    \[
    \bigcup_{n=1}^{\infty}E_n
    \subset
    \bigcup_{n=1}^{\infty}\bigcup_{k=1}^{\infty}A_{n,k}
    \]
    だから
    \begin{align*}
    \mu^*\left(\bigcup_{n=1}^{\infty}E_n\right)
    &\le
    \sum_{n=1}^{\infty}\sum_{k=1}^{\infty}\mu_0(A_{n,k}) \\
    &<
    \sum_{n=1}^{\infty}
    \left(
    \mu^*(E_n)+\frac{\eps}{2^n}
    \right) \\
    &=
    \sum_{n=1}^{\infty}\mu^*(E_n)+\eps.
    \end{align*}
    $\eps>0$ は任意だから
    \[
    \mu^*\left(\bigcup_{n=1}^{\infty}E_n\right)
    \le
    \sum_{n=1}^{\infty}\mu^*(E_n)
    \]
    である．
\end{enumerate}
\end{proof}

\begin{proposition}[前測度の拡張]
任意の $A \in \mathcal A$ に対して
\[
\mu^*(A)=\mu_0(A)
\]
が成り立つ．
\end{proposition}
\begin{proof}
$A \in \mathcal A$ とする．
$A$ 自身による1項被覆を使えば
\[
\mu^*(A)\le \mu_0(A)
\]
である．

逆向きの不等式を示す．
$A \subset \bigcup_{n=1}^{\infty}A_n$ を満たす任意の被覆 $\{A_n\}\subset\mathcal A$ を取る．
各 $n$ に対して
\[
B_1:=A\cap A_1,\qquad
B_n:=A\cap\left(A_n\setminus\bigcup_{k=1}^{n-1}A_k\right)\quad(n\ge 2)
\]
とおく．
有限加法族の閉性より各 $B_n$ は $\mathcal A$ に属する．
また $\{B_n\}$ は互いに素で，
\[
A=\bigsqcup_{n=1}^{\infty}B_n
\]
である．
しかも各 $B_n \subset A_n$ だから単調性より
\[
\mu_0(B_n)\le \mu_0(A_n)
\]
である．
前測度の可算加法性より
\[
\mu_0(A)
=
\sum_{n=1}^{\infty}\mu_0(B_n)
\le
\sum_{n=1}^{\infty}\mu_0(A_n).
\]
被覆の取り方について下限を取れば
\[
\mu_0(A)\le \mu^*(A)
\]
を得る．
\end{proof}

\begin{corollary}[外測度から測度へ]
Carathéodoryの定理により，
\[
\mu:=\mu^*|_{\mathcal M_{\mu^*}}
\]
は測度であり，
しかも
\[
\mu(A)=\mu_0(A)\qquad (A\in\mathcal A)
\]
である．
\end{corollary}
\begin{proof}
Carathéodoryの定理と直前の命題を合わせればよい．
\end{proof}

\begin{definition}[$\sigma$-有限]
有限加法族上の前測度 $\mu_0$，または $\sigma$-加法族上の測度 $\mu$ が
$\sigma$-有限であるとは，ある増大列
\[
X_1 \subset X_2 \subset \cdots
\]
で
\[
X=\bigcup_{n=1}^{\infty}X_n
\]
かつ
\[
\mu_0(X_n)<\infty
\qquad
\text{または}
\qquad
\mu(X_n)<\infty
\]
を満たすものが存在することをいう．
\end{definition}

\begin{theorem}[Carathéodory-Hahnの拡張定理]
$\mathcal S$ を集合 $X$ 上の半加法族，
$\mu_0$ をその上の前測度とする．
このとき，$\sigma(\mathcal S)$ 上の測度 $\mu$ で
\[
\mu(A)=\mu_0(A)\qquad (A\in\mathcal S)
\]
を満たすものが存在する．

さらに，$\overline{\mu}_0$ が $\mathcal A(\mathcal S)$ 上で
$\sigma$-有限ならば，この拡張は一意的である．
\end{theorem}
\begin{proof}
存在を示す．
$\mathcal A:=\mathcal A(\mathcal S)$ とおく．
先の命題により $\mathcal A$ は有限加法族であり，
$\overline{\mu}_0$ はその上の前測度である．
そこで $\overline{\mu}_0$ から外測度 $\mu^*$ を作ると，
Carathéodoryの定理により
\[
\widetilde{\mu}:=\mu^*|_{\mathcal M_{\mu^*}}
\]
は $\mathcal M_{\mu^*}$ 上の測度であり，
\[
\widetilde{\mu}(A)=\overline{\mu}_0(A)
\qquad
(A\in\mathcal A)
\]
である．

ここで $\mathcal M_{\mu^*}$ は $\sigma$-加法族で
$\mathcal A$ を含むから，
\[
\sigma(\mathcal S)=\sigma(\mathcal A)\subset \mathcal M_{\mu^*}
\]
である．
したがって
\[
\mu:=\widetilde{\mu}|_{\sigma(\mathcal S)}
\]
と定めれば，これは $\sigma(\mathcal S)$ 上の測度である．
また $\mathcal S \subset \mathcal A$ だから，
$A \in \mathcal S$ に対して
\[
\mu(A)=\widetilde{\mu}(A)=\overline{\mu}_0(A)=\mu_0(A)
\]
が成り立つ．

次に一意性を示す．
$\overline{\mu}_0$ が $\sigma$-有限であるとし，
$\sigma(\mathcal S)$ 上の別の測度 $\nu$ で
\[
\nu(A)=\mu_0(A)\qquad (A\in\mathcal S)
\]
を満たすものを取る．
$\mathcal A$ 上では $\nu$ と $\overline{\mu}_0$ は一致する．
実際，
\[
E=\bigsqcup_{j=1}^N A_j
\qquad
(A_j \in \mathcal S)
\]
ならば，$\nu$ の有限加法性より
\[
\nu(E)=\sum_{j=1}^N \nu(A_j)=\sum_{j=1}^N \mu_0(A_j)=\overline{\mu}_0(E)
\]
である．

$\sigma$-有限性より，
ある増大列 $\{X_n\}_{n=1}^{\infty}\subset\mathcal A$ が存在して
\[
X=\bigcup_{n=1}^{\infty}X_n,
\qquad
\overline{\mu}_0(X_n)<\infty
\]
を満たす．
任意の $E \in \sigma(\mathcal S)$ を取る．
各 $n$ に対して
\[
E_n:=E\cap X_n
\]
とおく．

まず
\[
\nu(E_n)\le \mu^*(E_n)
\]
を示す．
$E_n \subset \bigcup_{k=1}^{\infty}A_k$ を満たす任意の
$A_k \in \mathcal A$ を取ると，$\nu$ の可算劣加法性より
\[
\nu(E_n)\le \sum_{k=1}^{\infty}\nu(A_k)
=
\sum_{k=1}^{\infty}\overline{\mu}_0(A_k)
\]
である．
被覆について下限を取れば
\[
\nu(E_n)\le \mu^*(E_n)
\]
を得る．

逆向きの不等式を示す．
$X_n \setminus E$ を覆う任意の $A_k \in \mathcal A$ に対して
\[
X_n \subset E_n \cup \bigcup_{k=1}^{\infty}A_k
\]
だから
\[
\overline{\mu}_0(X_n)=\nu(X_n)\le \nu(E_n)+\sum_{k=1}^{\infty}\overline{\mu}_0(A_k).
\]
したがって
\[
\overline{\mu}_0(X_n)-\nu(E_n)\le \mu^*(X_n \setminus E).
\]
一方，$E_n \in \sigma(\mathcal S)\subset \mathcal M_{\mu^*}$ であり，
$X_n \in \mathcal A \subset \mathcal M_{\mu^*}$ だから
\[
\mu^*(X_n)=\mu^*(E_n)+\mu^*(X_n \setminus E)
\]
が成り立つ．
しかも $X_n \in \mathcal A$ なので
\[
\mu^*(X_n)=\overline{\mu}_0(X_n)
\]
である．
よって
\[
\mu^*(E_n)=\overline{\mu}_0(X_n)-\mu^*(X_n \setminus E)\le \nu(E_n)
\]
を得る．
以上より
\[
\nu(E_n)=\mu^*(E_n)
\qquad
(n \in \NN)
\]
である．

最後に $n\to\infty$ とする．
$E_n \uparrow E$ だから，$\nu$ の連続性より
\[
\nu(E)=\lim_{n\to\infty}\nu(E_n)
\]
である．
また $E \in \sigma(\mathcal S)\subset \mathcal M_{\mu^*}$ であるから，
$\mu^*|_{\mathcal M_{\mu^*}}$ は測度であり，
\[
\mu^*(E)=\lim_{n\to\infty}\mu^*(E_n)
\]
が成り立つ．
したがって
\[
\nu(E)
=
\lim_{n\to\infty}\nu(E_n)
=
\lim_{n\to\infty}\mu^*(E_n)
=
\mu^*(E)
=
\mu(E)
\]
である．
任意の $E \in \sigma(\mathcal S)$ に対して成り立つから，
$\nu=\mu$ である．
\end{proof}

\begin{remark}
存在の部分を Carathéodory-Hahn の拡張定理，
$\sigma$-有限性のもとでの一意性まで含めて Hopf の拡張定理と呼ぶことがある．
重要なのは，
\[
\text{半加法族}
\longrightarrow
\text{有限加法族}
\longrightarrow
\text{外測度}
\longrightarrow
\sigma\text{-加法族上の測度}
\]
という段階的な構成が一般に通用することである．
\end{remark}

\begin{example}[Lebesgue測度の構成]
$X=\RR^d$ とし，
\[
\mathcal S
:=
\left\{
\prod_{k=1}^d [a_k,b_k)
\;\middle|\;
a_k \le b_k
\right\}
\]
を半開区間全体の半加法族とする．
各 $I \in \mathcal S$ に対して
\[
\mu_0(I):=|I|
\]
と定める．

まず，$\mu_0$ が $\mathcal S$ 上の前測度であることを示す．
\[
I=\bigsqcup_{n=1}^{\infty}I_n
\qquad
(I,I_n \in \mathcal S)
\]
とする．
各 $N$ に対して
\[
\bigsqcup_{n=1}^N I_n \subset I
\]
だから，第2章の基本集合のJordan測度の単調性より
\[
\sum_{n=1}^N |I_n|
=
m_J\left(\bigsqcup_{n=1}^N I_n\right)
\le
|I|
\]
である．
$N\to\infty$ として
\[
\sum_{n=1}^{\infty}|I_n|\le |I|
\]
を得る．

逆向きの不等式は，第3章で示した
\[
\lambda^*(I)=|I|
\]
を用いれば
\[
|I|
=
\lambda^*(I)
\le
\sum_{n=1}^{\infty}\lambda^*(I_n)
=
\sum_{n=1}^{\infty}|I_n|
\]
と従う．
したがって
\[
|I|=\sum_{n=1}^{\infty}|I_n|
\]
であり，$\mu_0$ は確かに前測度である．

さらに $\mathcal A(\mathcal S)$ は，第2章で導入した基本集合全体である．
また $\overline{\mu}_0$ から生成される外測度は，第3章の
Lebesgue外測度 $\lambda^*$ と一致する．
実際，第3章では有界区間による可算被覆の下限を取ったが，
任意の有界区間は有限個の半開区間に分割でき，
体積の総和は変わらないからである．
よって，半開区間だけを使って被覆しても下限の値は変わらない．

したがって，この一般構成は
\[
\text{半開区間}
\longrightarrow
\text{基本集合}
\longrightarrow
\text{Lebesgue外測度}
\longrightarrow
\text{Lebesgue測度}
\]
という第2章から第5章までの流れを抽象化したものになっている．
また $\sigma(\mathcal S)=\mathcal B(\RR^d)$ であるから，
Carathéodory-Hahnの定理はまずBorel測度を与え，
その完備化としてLebesgue測度が得られる．
\end{example}

\begin{example}[数え上げ測度]
任意の集合 $X$ に対して
\[
\#(A):=
\begin{cases}
\text{$A$ の元の個数} & (A \text{ が有限集合のとき}),\\
\infty & (A \text{ が無限集合のとき})
\end{cases}
\qquad
(A \subset X)
\]
と定める．
これは $\mathcal P(X)$ 上の測度であり，数え上げ測度と呼ばれる．
\end{example}
\begin{proof}
$\#(\emptyset)=0$ は明らかである．
互いに素な列 $\{A_n\}_{n=1}^{\infty}\subset\mathcal P(X)$ を取る．

\begin{enumerate}
    \item ある $A_n$ が無限集合ならば
    \[
    \#\left(\bigcup_{n=1}^{\infty}A_n\right)=\infty
    \]
    であり，右辺も $\infty$ を含むから
    \[
    \sum_{n=1}^{\infty}\#(A_n)=\infty
    \]
    である．

    \item すべての $A_n$ が有限集合で，しかも
    非空なものが有限個しかなければ，
    通常の有限個の個数の加法性より
    \[
    \#\left(\bigcup_{n=1}^{\infty}A_n\right)
    =
    \sum_{n=1}^{\infty}\#(A_n)
    \]
    である．

    \item すべての $A_n$ が有限集合で，非空なものが無限個あるならば，
    和集合は無限集合であるから左辺は $\infty$ である．
    一方，各非空な $A_n$ について $\#(A_n)\ge 1$ だから
    \[
    \sum_{n=1}^{\infty}\#(A_n)=\infty
    \]
    である．
\end{enumerate}

以上により可算加法性が従う．
\end{proof}

\begin{example}[Dirac測度]
集合 $X$ と点 $x_0 \in X$ を取る．
$\mathcal P(X)$ 上の写像
\[
\delta_{x_0}(A)
:=
\begin{cases}
1 & (x_0 \in A),\\
0 & (x_0 \notin A)
\end{cases}
\qquad
(A \subset X)
\]
を Dirac 測度という．
\end{example}
\begin{proof}
$\delta_{x_0}(\emptyset)=0$ は明らかである．
互いに素な列 $\{A_n\}_{n=1}^{\infty}$ を取る．
$x_0$ を含む $A_n$ は高々1つしかないから，
\[
\delta_{x_0}\left(\bigcup_{n=1}^{\infty}A_n\right)
=
\sum_{n=1}^{\infty}\delta_{x_0}(A_n)
\]
が成り立つ．
したがって $\delta_{x_0}$ は測度である．
\end{proof}

\begin{remark}
Lebesgue測度は面積や体積を表す測度であるが，
測度という概念そのものはもっと広い．
数え上げ測度は「点の個数」を測り，
Dirac測度は「ある1点にだけ質量が集中している」ことを表す．
測度論の一般構成は，このように性質の異なる量を
同じ言葉で統一的に扱えることに本質がある．
\end{remark}

\section{解釈}

Carathéodory流の測度論は，
「面積を与える」という幾何学的問題を，
外測度と可測性という抽象的な言葉に翻訳したものである．
その利点は，
\begin{enumerate}
    \item 位相に依存しない
    \item 一般の集合 $X$ で使える
    \item 測度の存在を系統的に証明できる
\end{enumerate}
という点にある．

したがって，
Lebesgue流の
\[
\text{外測度}=\text{内測度}
\]
という定義は幾何学的直観を与え，
Carathéodory流の定義はその直観を抽象理論として完成させる役割を果たしている．
